\section{דיון}
תוצאות הפרויקט מצביעות על כך שניתן לפתח מודל חיזוי מבוסס נתונים להערכת משתני בית השורשים בחממה, גם כאשר מדידות היעד של pH ו-EC דלילות יחסית. שילוב נתוני אקלים, מיקרו-אקלים, השקיה, דישון, מצב הצמח ומדידות עבר אפשר למודל ללמוד את הדינמיקה הכללית של המצע ולספק תחזית עתידית בעלת ערך תפעולי.

\subsection{משמעות תוצאות מודל המיקרו-אקלים}
מודל המיקרו-אקלים הציג ביצועים טובים בחיזוי המשתנים המרכזיים בתוך החממה, ובהם טמפרטורת אוויר פנימית, לחות יחסית, קרינה פנימית, ET0 וטמפרטורת קרקע. תוצאות אלו חשובות משום שבית השורשים אינו מושפע ישירות רק מנתוני התחנה המטאורולוגית החיצונית, אלא מהתנאים בפועל בתוך סביבת הגידול.

הדיוק הגבוה של מודל המיקרו-אקלים מאפשר להשתמש בו כשכבת קלט יציבה עבור מודל בית השורשים. במילים אחרות, המודל יוצר בסיס סביבתי רציף שמאפשר לחזות את תגובת המצע גם כאשר אין מדידה ישירה מלאה של כל משתני החממה בכל רגע.

\subsection{משמעות תוצאות מודל בית השורשים}
מודל בית השורשים הצליח לחזות את ערכי pH ו-EC ברמת דיוק טובה ביחס למודל נאיבי המבוסס על שמירת הערך האחרון שנמדד. עבור pH התקבלה יציבות גבוהה יחסית הן במבחן Walk-Forward והן במבחן Holdout, דבר המצביע על כך שהמודל מצליח ללמוד את דפוסי השינוי המרכזיים של החומציות במצע.

עבור EC התקבלו גם כן ביצועים טובים מבחינת MAE ו-\Rtwo{}, אך החיזוי נותר מורכב יותר. הסיבה המרכזית לכך היא שערכי EC נמוכים גורמים לכך שגם סטייה אבסולוטית קטנה מתורגמת לשגיאה יחסית גבוהה. לכן, בהערכת ביצועי EC חשוב להתייחס לא רק לאחוזי השגיאה, אלא גם למשמעות האגרונומית של הסטייה.

\subsection{תרומת הגישה ההיברידית}
אחד היתרונות המרכזיים של המודל הוא השילוב בין רכיב \GB{} לבין רכיב ליניארי חסין. רכיב ה־\GB{} מאפשר ללמוד קשרים מורכבים ולא-ליניאריים מתוך דפוסים שנצפו בנתוני העבר, ואילו הרכיב הליניארי מסייע בהתמודדות עם מגמות ואירועים תפעוליים בעלי השפעה גבוהה.

שילוב זה מאפשר למודל לבצע אינטרפולציה בתוך תחומי התנהגות מוכרים, ובמקביל לבצע אקסטרפולציה זהירה במצבים שבהם מופיעים אירועים פחות שכיחים. בכך המודל מתאים יותר למערכת חקלאית דינמית, שבה קיימים גם דפוסים מחזוריים וגם אירועים נקודתיים חריגים.

\subsection{מגבלות המחקר}
למרות התוצאות החיוביות, קיימות מספר מגבלות שיש לקחת בחשבון. ראשית, מספר מדידות היעד של pH ו-EC היה מוגבל יחסית, ולכן יש צורך בהמשך איסוף נתונים כדי לשפר את הכללת המודל. שנית, המודל נבחן על בסיס עונת גידול אחת, ולכן עדיין לא ניתן לקבוע בוודאות את ביצועיו בעונות אחרות, בגידולים אחרים או בחממות שונות.

בנוסף, בדיקת ה-Holdout בוצעה מתוך אותה עונת גידול, ולכן היא אינה מחליפה בדיקה מלאה על עונת גידול עתידית עצמאית. כמו כן, המדידות הידניות עצמן עשויות לכלול רעש, שונות מרחבית או טעויות מדידה, ולכן חלק מהשגיאה של המודל עשוי לנבוע גם מאי-ודאות במדידת היעד.

\subsection{משמעות יישומית}
המשמעות היישומית של הפרויקט היא שניתן להשתמש במודל כ־\SoftSensor{} תומך החלטה לניהול בית השורשים. המודל אינו מחליף מדידות ידניות או שיקול דעת אגרונומי, אך הוא יכול לסייע בהערכת המגמה בין מועדי המדידה, בזיהוי מוקדם של סיכון לשינויי pH או EC, ובתכנון מושכל יותר של השקיה ודישון.

בכך הפרויקט מדגים מעבר מגישה תגובתית, שבה מתקנים חריגות לאחר שהן נמדדו, לגישה פרואקטיבית יותר, שבה ניתן להעריך מראש את כיוון השינוי הצפוי ולתמוך בקבלת החלטות בזמן.
